\chapter{Diagramas de arquitectura}

\section{Notación y fuente de los diagramas}
Los diagramas se modelaron en \textbf{ArchiMate} con la herramienta Archi, en un solo modelo versionado
(\cod{docs/diagramas/safemode-arquitectura.archimate}); las figuras de este capítulo son las vistas 4 a 16 de ese modelo, exportadas con Archi
en línea de comandos (\cod{docs/arquitectura/exportar-vistas.py}). Describen \textbf{lo implementado}: cada puerto, ruta, clase, tabla y servicio
se verificó contra el código. Las tres vistas anteriores del equipo (1 a 3, de la etapa de diseño) siguen en el modelo como historia.

ArchiMate no trae la notación UML, así que cada tipo de diagrama que pide el requerimiento se expresa con los elementos que el estándar sí
tiene. La Tabla~\ref{tab:uml} dice cuál vista cubre cada uno y cómo se lee.

\begin{table}[H]
\centering\small
\caption{Diagramas requeridos y la vista ArchiMate que los cubre.}\label{tab:uml}
\begin{tabularx}{\textwidth}{lllY}
\toprule
\textbf{Diagrama} & \textbf{Vista} & \textbf{Figura} & \textbf{Cómo se expresa en ArchiMate} \\
\midrule
Contexto & 4 & \ref{fig:contexto} & El sistema como un componente de aplicación rodeado de actores y sistemas externos (en gris) \\
Componentes & 5 y 6 & \ref{fig:comp-flujo}, \ref{fig:comp-puesto} & Componentes de aplicación con composición (partes de un proceso) y flujos entre ellos \\
Clases & 7, 8 y 9 & \ref{fig:clases-1}--\ref{fig:clases-3} & Cada clase es un componente; \emph{Realization} = implementa la interfaz; \emph{Composition} = contiene; \emph{Serving} = presta servicio a \\
Entidad-relación & 10 & \ref{fig:er} & Objetos de datos (tablas, documentos, claves) con asociaciones N:1 \\
Secuencia & 11 y 12 & \ref{fig:secuencia-alerta}, \ref{fig:secuencia-login} & Vista dinámica: componentes y flujos numerados en el orden de la interacción \\
Despliegue & 13 y 14 & \ref{fig:despliegue-actual}, \ref{fig:despliegue-objetivo} & Nodos, software de sistema y redes de la capa de tecnología \\
Actividades & 15 y 16 & \ref{fig:act-1}, \ref{fig:act-2} & Procesos con relaciones de disparo, eventos y uniones (O = decisión; Y = unión) \\
\bottomrule
\end{tabularx}
\end{table}

\textbf{Convenciones.} Gris = sistema externo o fuera de nuestro control; azul = capa de aplicación; verde = tecnología; amarillo =
negocio; naranja (Figura~\ref{fig:despliegue-objetivo}) = diseñado y calculado pero \emph{no desplegado}. Las líneas punteadas con flecha son
flujos; las continuas con rombo, composición.

\section{Contexto}
El sistema se ve como una sola pieza rodeada de quienes lo usan y de los sistemas con los que habla. El personal de seguridad y el administrador
reciben servicio del sistema; las cámaras se ven en el software de cámaras (VMS), cuya ventana compartida es la entrada de imagen; Google
verifica la identidad; Telegram recibe avisos de texto; el servicio de visión de NVIDIA es opcional; y GitHub Actions y SonarCloud
despliegan y analizan el código. El video no se guarda en SafeMode.

\vista{04-contexto.png}{Contexto del sistema (vista 4).}{fig:contexto}

\section{Componentes}
La Figura~\ref{fig:comp-flujo} sigue el camino de una alerta: el panel recibe la ventana del VMS y envía un fotograma por cámara al
\emph{Camera Feed}; el detector produce detecciones que el seguimiento convierte en eventos (y fotos); el motor de alertas, dentro del gateway,
lee los retiros nuevos, crea la alerta y la entrega al servidor WebSocket y al servicio de avisos. El gateway tiene seis partes (composición) y
lanza el tracker con un comando fijo. La Figura~\ref{fig:comp-puesto} muestra el Release~2: el puesto compartido (zonas, historial y presencia) en
\emph{collaboration-service}, los permisos en \emph{access-service} (que la colaboración consulta antes de cada escritura y que falla cerrada), los
tres motores de datos y la observabilidad.

\vista{05-componentes-flujo.png}{Componentes: de la cámara a la alerta (vista 5).}{fig:comp-flujo}
\vista{06-componentes-puesto.png}{Componentes: puesto compartido, permisos y observabilidad (vista 6).}{fig:comp-puesto}

\section{Clases}
Se muestran las clases que sostienen las decisiones de diseño, en tres vistas. La primera (Figura~\ref{fig:clases-1}): el
\cod{ObjectTracker} (uno por cámara) contiene objetos y personas seguidos, usa descriptores de apariencia y una huella de la zona para
verificar un retiro, y delega la descripción por IA en una interfaz que implementa el cliente de NVIDIA. La segunda (Figura~\ref{fig:clases-2}):
los almacenes de eventos y alertas comparten conexión con reconexión automática y dialecto (H2 o PostgreSQL), y el motor de alertas combina
eventos y reglas para crear alertas únicas que avisa a sus oyentes. La tercera (Figura~\ref{fig:clases-3}): el \cod{GatewayRunner} arma el
servidor HTTP, el WebSocket, la autenticación y el servicio de avisos; la sesión y la verificación del token son interfaces con dos
implementaciones (memoria o base) y una (Google), lo que permitió replicar el gateway sin cambiar su código.

\vista{07-clases-deteccion.png}{Clases (1/3): detección y seguimiento (vista 7).}{fig:clases-1}
\vista{08-clases-persistencia.png}{Clases (2/3): persistencia y motor de alertas (vista 8).}{fig:clases-2}
\vista{09-clases-gateway.png}{Clases (3/3): gateway, sesiones y avisos (vista 9).}{fig:clases-3}

\section{Entidad-relación}
El modelo de datos tiene cuatro grupos. En la base de eventos (H2 por defecto, PostgreSQL con una URL): \cod{object_presence_event} y
\cod{alert} (una alerta por evento, garantizada por \cod{UNIQUE}), \cod{camera} y \cod{operator_session} (la huella SHA-256 de la cookie, nunca la
cookie). En PostgreSQL de acceso: puestos, operadores, membresías (rol y permiso de configurar) y auditoría de concesiones y revocaciones. En MongoDB:
el documento del puesto con su versión y el historial de cambios (índice único por puesto y versión, retención de 90 días). En Redis: presencia y
caché de permisos, con vencimiento.

\vista{10-entidad-relacion.png}{Modelo de datos: tablas, documentos y claves (vista 10).}{fig:er}

\section{Secuencia}
La primera secuencia (Figura~\ref{fig:secuencia-alerta}) es la promesa de valor: del retiro de un objeto al aviso en el panel en menos de
5~s, con el presupuesto de tiempo anotado. La segunda (Figura~\ref{fig:secuencia-login}) es el inicio de sesión con Google: el ID token se
verifica en el servidor (se le pregunta a Google y se comprueban emisor, audiencia y vigencia), el correo debe estar en la lista blanca, la sesión se guarda por su huella y el
navegador recibe una cookie \cod{HttpOnly}; el WebSocket revalida la sesión.

\vista{11-secuencia-alerta.png}{Secuencia: del retiro de un objeto al aviso en el panel (vista 11).}{fig:secuencia-alerta}
\vista{12-secuencia-login.png}{Secuencia: inicio de sesión con Google (vista 12).}{fig:secuencia-login}

\section{Despliegue}
La Figura~\ref{fig:despliegue-actual} es el despliegue público de hoy: una instancia EC2 que ejecuta Docker Compose con Caddy delante (TLS,
\cod{forward_auth} para \cod{/frames} y \cod{/grafana}) y todos los servicios detrás; un solo servidor, por lo que su caída es total. La
Figura~\ref{fig:despliegue-objetivo} pone lado a lado el \textbf{prototipo validado} (dos réplicas detrás de un balanceador con PostgreSQL, medido en
la sección~\ref{sec:disponibilidad}) y el \textbf{diseño objetivo} en AWS (Route~53, balanceador multi-AZ, dos instancias en zonas distintas y RDS
Multi-AZ), señalado como no desplegado.

\vista{13-despliegue-actual.png}{Despliegue actual: servidor único en AWS con Docker Compose (vista 13).}{fig:despliegue-actual}
\vista{14-despliegue-objetivo.png}{Prototipo de escalado validado y diseño objetivo en AWS (vista 14).}{fig:despliegue-objetivo}

\section{Actividades}
El flujo se parte en dos. En la Figura~\ref{fig:act-1}, cada fotograma se detecta y se empareja con lo ya seguido; un objeto se registra cuando lleva 3~s
quieto; y se da por retirado si desaparece 3 fotogramas o se desplaza más de 25~px en 2 seguidos, siempre que la verificación por imagen
(huella de 24$\times$24 puntos, sin contar a las personas que tapan) confirme que cambió. En la Figura~\ref{fig:act-2}, el retiro se clasifica
según quién estaba cerca y se traduce a una alerta de severidad alta, media o baja; la alerta se crea una sola vez, se publica a todos los
operadores y, si es alta o media, sale el aviso de texto.

\vista{15-actividades-retiro.png}{Actividades (1/2): del fotograma al retiro confirmado (vista 15).}{fig:act-1}
\vista{16-actividades-alerta.png}{Actividades (2/2): del retiro a la alerta y el aviso (vista 16).}{fig:act-2}
